\section{Annex B: Dense, Recurrent, and Memory-Augmented Transformers}
\label{sec:annex-transformer}

This comprehensive annex synthesizes modern transformer architectures beyond the standard softmax baseline. The field has evolved toward six primary paradigms: dense global attention with variants, key-value head compression via grouped-query attention, recurrent state compression with decay, selective state-space models, continuous-depth numerical integration, memory-augmented architectures, and hybrid approaches. Understanding this taxonomy illuminates fundamental tradeoffs between expressiveness, computational cost, memory footprint, and deployment simplicity.

\subsection{Dense Attention Baselines: Standard and Sigmoid}

\subsubsection{Standard Softmax Attention}
Standard multi-head self-attention \citep{vaswani_attention_2017} computes scaled dot-product similarity between token embeddings:

\begin{algorithmic}[1]
\State \textbf{Input:} Query matrix $\mathbf{Q} \in \mathbb{R}^{n \times d}$, Key matrix $\mathbf{K} \in \mathbb{R}^{n \times d}$, Value matrix $\mathbf{V} \in \mathbb{R}^{n \times d}$
\State $\text{scores} \gets \frac{\mathbf{Q}\mathbf{K}^\top}{\sqrt{d}} \in \mathbb{R}^{n \times n}$ \Comment{compute pairwise similarities}
\State $\text{attention\_weights} \gets \texttt{softmax}(\text{scores}, \text{axis}=1) \in \mathbb{R}^{n \times n}$ \Comment{normalize across keys}
\State \textbf{Output:} $\mathbf{Y} = \text{attention\_weights} \cdot \mathbf{V} \in \mathbb{R}^{n \times d}$ \Comment{weighted value combination}
\end{algorithmic}

For autoregressive generation, KV caching stores past keys and values to avoid $\mathcal{O}(n^2)$ recomputation. However, this linearly growing cache occupies $\sim n \cdot d_{\text{hidden}}$ bytes, which can consume hundreds of gigabytes for billion-parameter models. Standard attention achieves \textbf{perfect expressiveness} within the context window---any token can attend to any other token with learned weights---but pays the price of dense computation.

\begin{itemize}[leftmargin=*]
    \item \textbf{Training complexity}: $\mathcal{O}(n^2 \cdot d)$ time, $\mathcal{O}(n^2)$ space (attention matrix materialization).
    \item \textbf{Inference complexity}: $\mathcal{O}(n)$ time per token, $\mathcal{O}(n)$ space (KV cache).
    \item \textbf{Strengths}: Unparalleled expressiveness; perfect history recall; highly parallelizable.
    \item \textbf{Weaknesses}: Quadratic bottleneck prohibits very long contexts; KV cache dominates memory during generation.
\end{itemize}

\subsubsection{Sigmoid Attention}
Sigmoid attention replaces row-wise softmax with element-wise sigmoid activation \citep{ramapuram_theory_2024}:

\begin{algorithmic}[1]
\State \textbf{Input:} $\mathbf{Q}, \mathbf{K}, \mathbf{V}$ as above, learnable bias $\mathbf{b} \in \mathbb{R}^{n \times n}$
\State $\text{logits} \gets \frac{\mathbf{Q}\mathbf{K}^\top}{\sqrt{d}} + \mathbf{b}$
\State $\text{attention\_weights} \gets \sigma(\text{logits})$ \Comment{element-wise sigmoid, not softmax}
\State \textbf{Output:} $\mathbf{Y} = \text{attention\_weights} \odot \mathbf{V}$
\end{algorithmic}

Unlike softmax, sigmoid does not enforce a probability distribution (weights need not sum to 1), enabling stronger token independence. Theoretical analysis via mixture-of-experts shows sigmoid achieves superior sample complexity: $\mathcal{O}(n^{-0.51})$ convergence for ReLU experts versus softmax's $\mathcal{O}(n^{-0.24})$. However, empirical training revealed gradient instabilities at scale. The remedy is \textbf{hybrid-norm}---adding normalization after the attention output---which stabilizes gradients without sacrificing the theoretical benefits of element-wise gating.

\begin{itemize}[leftmargin=*]
    \item \textbf{Training complexity}: $\mathcal{O}(n^2 \cdot d)$ (identical to standard), but element-wise operations enable 17\% inference speedup via FlashSigmoid.
    \item \textbf{Inference complexity}: $\mathcal{O}(n)$ per token with KV cache (asymptotically same, but lower constant factors).
    \item \textbf{Strengths}: Overcomes zero-sum competition; avoids row-wise synchronization; hardware-friendly implementation.
    \item \textbf{Weaknesses}: Requires careful stabilization (hybrid-norm); training instability at large scales without auxiliary loss.
\end{itemize}

\subsubsection{Grouped-Query Attention (GQA)}
Grouped-query attention \citep{ainslie_gqa_2023} bridges multi-head attention (MHA) and multi-query attention (MQA) by partitioning query heads into groups, each sharing a single key-value head. If $h$ is the total number of attention heads, $h_k$ is the number of key-value heads, and $G = h / h_k$ is the group size, then every group of $G$ query heads shares one KV head:

\begin{algorithmic}[1]
\State \textbf{Input:} $\mathbf{Q} \in \mathbb{R}^{n \times h \cdot d_h}$, $\mathbf{K} \in \mathbb{R}^{n \times h_k \cdot d_h}$, $\mathbf{V} \in \mathbb{R}^{n \times h_k \cdot d_h}$
\State $\mathbf{K}_{\text{full}} \gets \texttt{repeat\_interleave}(\mathbf{K}, G)$ \Comment{expand from $h_k$ to $h$ heads}
\State $\mathbf{V}_{\text{full}} \gets \texttt{repeat\_interleave}(\mathbf{V}, G)$
\State $\text{scores} \gets \frac{\mathbf{Q}\mathbf{K}_{\text{full}}^\top}{\sqrt{d_h}} \in \mathbb{R}^{n \times n}$ \Comment{standard dot product after expansion}
\State $\text{attention\_weights} \gets \texttt{softmax}(\text{scores}, \text{axis}=1)$
\State \textbf{Output:} $\mathbf{Y} = \text{attention\_weights} \cdot \mathbf{V}_{\text{full}} \in \mathbb{R}^{n \times h \cdot d_h}$
\end{algorithmic}

GQA interpolates smoothly between MHA ($h_k = h$, $G = 1$) and MQA ($h_k = 1$, $G = h$). The key advantage is a tunable tradeoff between KV-cache footprint and representational capacity. Each additional KV head consumes $2 \cdot d_h \cdot n$ extra bytes in the cache, so reducing $h_k$ yields proportional memory savings.

\begin{itemize}[leftmargin=*]
    \item \textbf{Training complexity}: $\mathcal{O}(n^2 \cdot h \cdot d_h)$ (identical to MHA after head expansion).
    \item \textbf{Inference complexity}: $\mathcal{O}(n)$ per token, $\mathcal{O}(h_k \cdot n \cdot d_h)$ KV-cache space.
    \item \textbf{Strengths}: Flexible head-count tradeoff; proven at scale (Llama 2 \& 3 use $h_k = 8$ KV heads with $h = 32$ query heads); reduces KV-cache by $G \times$; supports fast inference with FlashAttention.
    \item \textbf{Weaknesses}: Fewer KV heads can reduce representational capacity in small models; requiring $h_k$ to divide $h$ constrains architectural choices.
\end{itemize}

\subsection{Recurrent and Retentive Architectures}

\subsubsection{Retentive Networks (RetNet)}
RetNet \citep{sun_retentive_2023} unifies three computation paradigms: parallel training, recurrent inference, and chunkwise deployment. Its core innovation is the \textbf{retention mechanism}, which uses a fixed exponential decay matrix to model temporal importance:

\begin{algorithmic}[1]
\State \textbf{Parallel (training) form:}
\State $\text{decay\_matrix}[i,j] \gets \gamma^{i-j}$ for $i \geq j$, else $0$ \Comment{causal exponential decay}
\State $\text{decay\_matrix}[i,j] \gets 0$ for $i < j$ \Comment{causal masking}
\State $\mathbf{Y}_{\text{parallel}} \gets (\mathbf{Q}\mathbf{K}^\top \odot \text{decay\_matrix}) \mathbf{V}$ \Comment{element-wise multiplication with decay}
\end{algorithmic}

\begin{algorithmic}[1]
\State \textbf{Recurrent (inference) form:}
\For{$t = 1, \ldots, n$}
    \State $\mathbf{s}_t \gets \gamma \mathbf{s}_{t-1} + \mathbf{k}_t \mathbf{v}_t^\top$ \Comment{state update with decay}
    \State $\mathbf{y}_t \gets \mathbf{q}_t \mathbf{s}_t$ \Comment{output via query-state interaction}
\EndFor
\end{algorithmic}

The decay scalar $\gamma \in (0,1)$ controls the temporal window. RetNet uses \textbf{multi-scale retention} with different $\gamma$ values per head (e.g., $\gamma = 1 - 2^{-5}, 1 - 2^{-6}, \ldots$), allowing short-term and long-term dependencies simultaneously. Chunkwise recurrent mode divides sequences into chunks, processes each chunk in parallel, and threads a recurrent state between chunks.

\begin{itemize}[leftmargin=*]
    \item \textbf{Training complexity}: $\mathcal{O}(n^2 \cdot d)$ (parallel), or $\mathcal{O}(n \cdot c \cdot d)$ (chunkwise recurrent with chunk size $c$).
    \item \textbf{Inference complexity}: $\mathcal{O}(1)$ per token, $\mathcal{O}(d^2)$ state space (fixed matrix).
    \item \textbf{Strengths}: Constant-time inference; triple computation paradigm; multi-scale consolidation.
    \item \textbf{Weaknesses}: Fixed decay imposes rigid inductive bias; may truncate learned long-range patterns.
\end{itemize}

\subsubsection{Mamba: Selective State-Space Models}
Mamba \citep{gu_mamba_2023} frames recurrence as a continuous-time dynamical system with input-dependent parameters, achieving both linear training and constant-time inference:

\begin{algorithmic}[1]
\State \textbf{Continuous dynamics:} $h'(t) = \mathbf{A} h(t) + \mathbf{B} x(t)$, \quad $y(t) = \mathbf{C} h(t)$
\State \textbf{Discretization with step size $\Delta_t$:}
\State $\bar{\mathbf{A}}_t \gets \exp(\Delta_t \mathbf{A})$
\State $\bar{\mathbf{B}}_t \gets (\Delta_t \mathbf{A})^{-1} (\exp(\Delta_t \mathbf{A}) - \mathbf{I}) \Delta_t \mathbf{B}_t$
\State \textbf{Recurrent update:}
\For{$t = 1, \ldots, n$}
    \State $\Delta_t \gets \text{softplus}(\text{Linear}(x_t))$ \Comment{input-dependent step size}
    \State $\mathbf{B}_t \gets \text{Linear}(x_t)$, \quad $\mathbf{C}_t \gets \text{Linear}(x_t)$ \Comment{input-dependent projection matrices}
    \State $h_t \gets \bar{\mathbf{A}}_t h_{t-1} + \bar{\mathbf{B}}_t x_t$ \Comment{state transition}
    \State $y_t \gets \mathbf{C}_t h_t$ \Comment{output projectionContinuous}
\EndFor
\end{algorithmic}

The critical innovation is that $\mathbf{A}$ (the system matrix) is \textit{not} input-dependent, but $\Delta_t$, $\mathbf{B}_t$, and $\mathbf{C}_t$ are, making the system \textbf{time-varying}. This selectivity allows the model to \textit{ignore} irrelevant information by setting $\Delta_t$ near zero, effectively creating a gate. A hardware-aware parallel scan algorithm implements the recurrence efficiently on GPUs by fusing computation within SRAM, avoiding expensive HBM bandwidth.

\begin{itemize}[leftmargin=*]
    \item \textbf{Training complexity}: $\mathcal{O}(n \cdot d)$ via hardware-aware scan (linear in sequence length).
    \item \textbf{Inference complexity}: $\mathcal{O}(1)$ per token, $\mathcal{O}(d)$ state (hidden vector).
    \item \textbf{Strengths}: Achieves linear training and constant-inference simultaneously; practical on long sequences (millions of tokens).
    \item \textbf{Weaknesses}: State vector compression can weaken exact copying and dense associative recall versus full attention.
\end{itemize}

\subsection{Continuous-Depth Transformers: ODE Integration}

\subsubsection{ODE Transformer}
The ODE Transformer \citep{zhang_continuous_2021} interprets network depth as numerical integration of a continuous dynamical system, using higher-order Runge-Kutta solvers to reduce truncation error:

\begin{algorithmic}[1]
\State \textbf{Continuous formulation:} $\frac{dh(t)}{dt} = f_\theta(h(t), t)$ where $f_\theta$ is the transformer sub-network
\State \textbf{Runge-Kutta-4 discrete approximation:}
\State $k_1 \gets f_\theta(h_t, t)$
\State $k_2 \gets f_\theta(h_t + \tfrac{1}{2}k_1, t + \tfrac{1}{2}\Delta t)$
\State $k_3 \gets f_\theta(h_t + \tfrac{1}{2}k_2, t + \tfrac{1}{2}\Delta t)$
\State $k_4 \gets f_\theta(h_t + k_3, t + \Delta t)$
\State $h_{t+1} \gets h_t + \frac{\Delta t}{6}(k_1 + 2k_2 + 2k_3 + k_4)$ \Comment{weighted combination}
\end{algorithmic}

Instead of simple Euler residual connections $h_{t+1} = h_t + f_\theta(h_t)$, the RK4 block computes four intermediate evaluations and combines them with classical RK4 weights. To avoid vanishing gradients, the architecture introduces \textbf{learned gating} that interpolates between intermediate approximations:

\begin{algorithmic}[1]
\State $g \gets \sigma(\text{Linear}([k_1, k_2, k_3, k_4]))$ \Comment{learnable gate}
\State $h_{t+1} \gets h_t + g \cdot k_1 + (1-g) \cdot k_2$ \Comment{interpolated refinement}
\end{algorithmic}

This formulation reduces the effective number of parameters through weight sharing---the same $f_\theta$ is evaluated multiple times---while providing richer trajectory refinement.

\begin{itemize}[leftmargin=*]
    \item \textbf{Training complexity}: $\mathcal{O}(k \cdot n^2 \cdot d)$ where $k$ is the RK order (4 for RK4).
    \item \textbf{Inference complexity}: $\mathcal{O}(k \cdot n)$ per token (higher constant overhead).
    \item \textbf{Strengths}: Significantly higher accuracy on generation tasks (state-of-the-art BLEU); parameter-efficient via weight sharing.
    \item \textbf{Weaknesses}: High per-step compute; inference latency increases by constant factor $k$; complex gating required for stability.
\end{itemize}

\subsection{Test-Time Memory: Titans}

The Titans architecture \citep{behrouz_titans_2025} introduces an orthogonal dimension: instead of static weights, the model maintains a learnable memory that is \textit{updated during inference} based on a surprise-driven signal:

\begin{algorithmic}[1]
\State \textbf{Local short-term attention:} Apply standard or sparse attention within a fixed context window $c$
\State $y_t^{\text{local}} \gets \text{Attention}(q_t, k_{[t-c:t]}, v_{[t-c:t]})$
\State \textbf{Memory update signal (surprise):}
\State $S_t \gets \eta_t S_{t-1} - \theta_t \nabla_\ell(\mathcal{M}_{t-1}; x_t)$ \Comment{surprise as gradient of loss w.r.t. memory}
\State $\mathcal{M}_t \gets (1 - \alpha_t) \mathcal{M}_{t-1} + S_t$ \Comment{memory updated via EMA}
\State \textbf{Long-term memory retrieval:}
\State $y_t^{\text{memory}} \gets \mathcal{M}_t^* (q_t)$ \Comment{query memory module for retrieved information}
\State \textbf{Output combination:}
\State $y_t \gets y_t^{\text{local}} + \text{gate}(y_t^{\text{memory}})$ \Comment{combine local and retrieved memory}
\end{algorithmic}

The memory module $\mathcal{M}$ is literally updated during the forward pass by computing gradients of an associative loss and applying SGD steps with momentum. The decay rates $\eta_t, \theta_t, \alpha_t$ are themselves input-dependent, allowing the model to switch memory paradigms when context shifts.

\begin{itemize}[leftmargin=*]
    \item \textbf{Training complexity}: Roughly $\mathcal{O}(c^2 + n \cdot f)$ where $c$ is local window and $f$ is memory update overhead.
    \item \textbf{Inference complexity}: $\mathcal{O}(c^2)$ local attention plus $\mathcal{O}(1)$ memory retrieval per token.
    \item \textbf{Strengths}: Handles extreme context lengths; enables true associative recall; memory adapts to input.
    \item \textbf{Weaknesses}: Significantly more complex; inference includes gradient computation; higher coordination overhead.
\end{itemize}

\subsection{Multi-Head Latent Attention (MLA)}
Multi-Head Latent Attention (MLA) \citep{mehta_mla_2025} compresses the per-token key--value state into a single low-rank latent vector $c_{KV}$ of rank $r_{kv}$, which is the only quantity materialised in the KV cache during decoding. A pair of up-projection matrices reconstructs the full multi-head keys and values on the fly, while RoPE is applied to the decompressed query and key tensors so that positional structure is preserved. The authors show that the Pareto-optimal latent width sits near $r_{kv} = d_{\text{hidden}} / 2$, at which point MLA halves the KV cache yet matches multi-head attention (MHA) quality on small models.

\subsubsection{Mathematical Formulation}
\[
c_{KV} = W_{DKV}\, x \in \mathbb{R}^{r_{kv}}, \qquad
k = W_{UK}\, c_{KV}, \qquad
v = W_{UV}\, c_{KV}, \qquad
q = W_{Q}\, x.
\]
\[
\tilde q = \text{RoPE}(q), \quad \tilde k = \text{RoPE}(k), \qquad
\text{attn} = \texttt{softmax}\!\left(\frac{\tilde q\, \tilde k^\top}{\sqrt{d_h}}\right) v.
\]

\subsubsection{Algorithmic Pseudocode}
\begin{algorithmic}[1]
\State \textbf{Input:} hidden state $x \in \mathbb{R}^{n \times d}$, rank $r_{kv}$, projections $W_{DKV}, W_{UK}, W_{UV}, W_{Q}$
\State $c_{KV} \gets W_{DKV}\, x$ \Comment{down-project to latent of rank $r_{kv}$ (cached at inference)}
\State $k \gets W_{UK}\, c_{KV}$, \quad $v \gets W_{UV}\, c_{KV}$ \Comment{up-project to full heads}
\State $q \gets W_{Q}\, x$
\State $\tilde q \gets \text{RoPE}(q)$, \quad $\tilde k \gets \text{RoPE}(k)$ \Comment{apply rotary embeddings}
\State $\text{weights} \gets \texttt{softmax}(\tilde q\, \tilde k^\top / \sqrt{d_h})$
\State \textbf{Output:} $\mathbf{Y} = \text{weights} \cdot v$
\end{algorithmic}

\subsubsection{Key Characteristics}
\begin{itemize}[leftmargin=*]
    \item \textbf{Training complexity}: $\mathcal{O}(n^2 \cdot d)$ (full attention on decompressed heads).
    \item \textbf{Inference complexity}: $\mathcal{O}(n)$ per token, $\mathcal{O}(r_{kv} \cdot n)$ KV-cache space (latent only).
    \item \textbf{Trainable}: fully trainable; $r_{kv}$ is a hyperparameter (Pareto-optimal $\approx d/2$).
    \item \textbf{Strengths}: Halves KV cache versus MHA at matched quality; RoPE-compatible; clean latent bottleneck.
    \item \textbf{Limitations}: Up-projection adds decode FLOPs; latent rank is a fixed global budget shared across heads.
\end{itemize}

\subsection{Group-Query Latent Attention (GQLA)}
Group-Query Latent Attention (GQLA) \citep{meng_gqla_2026} is a minimal modification of MLA that exposes two algebraically equivalent decoding paths over the same trained weights: an MQA-absorb path (fusing the up-projection into the query, optimal on H100-class hardware) and a GQA-expanded path (materialising full keys/values, optimal on H20 / multi-token-prediction hardware). The two paths are mathematically identical because the up-projection is linear, so $\texttt{softmax}(q_{\text{absorbed}} \cdot c_{KV}^\top) = \texttt{softmax}((q \cdot W_{UK}^\top) \cdot c_{KV}^\top)$, allowing hardware-adaptive dispatch with no retraining.

\subsubsection{Mathematical Formulation}
Same latent as MLA:
\[
c_{KV} = W_{DKV}\, x, \qquad k = W_{UK}\, c_{KV}, \qquad v = W_{UV}\, c_{KV}, \qquad q = W_{Q}\, x.
\]
\[
\text{Path A (MQA-absorb): } \text{attn} = \texttt{softmax}\!\left(\frac{(q\, W_{UK}^\top)\, c_{KV}^\top}{\sqrt{d_h}}\right) (W_{UV}\, c_{KV}).
\]
\[
\text{Path B (GQA-expanded): } \text{attn} = \texttt{softmax}\!\left(\frac{q\, k^\top}{\sqrt{d_h}}\right) v, \quad k = W_{UK} c_{KV},\ v = W_{UV} c_{KV}.
\]

\subsubsection{Algorithmic Pseudocode}
\begin{algorithmic}[1]
\State \textbf{Input:} $x$, weights $W_{DKV}, W_{UK}, W_{UV}, W_{Q}$, hardware flag $\text{hw} \in \{\text{H100}, \text{H20}\}$
\State $c_{KV} \gets W_{DKV}\, x$ \Comment{shared latent, cached}
\State $q \gets W_{Q}\, x$
\If{$\text{hw} == \text{H100}$} \Comment{MQA-absorb path}
    \State $q_{\text{abs}} \gets q\, W_{UK}^\top$ \Comment{absorb up-projection into query}
    \State $\text{weights} \gets \texttt{softmax}(q_{\text{abs}}\, c_{KV}^\top / \sqrt{d_h})$
    \State $\mathbf{Y} \gets \text{weights} \cdot (W_{UV}\, c_{KV})$
\Else \Comment{GQA-expanded path}
    \State $k \gets W_{UK}\, c_{KV}$, \quad $v \gets W_{UV}\, c_{KV}$
    \State $\text{weights} \gets \texttt{softmax}(q\, k^\top / \sqrt{d_h})$
    \State $\mathbf{Y} \gets \text{weights} \cdot v$
\EndIf
\State \textbf{Output:} $\mathbf{Y}$
\end{algorithmic}

\subsubsection{Key Characteristics}
\begin{itemize}[leftmargin=*]
    \item \textbf{Training complexity}: $\mathcal{O}(n^2 \cdot d)$, identical to MLA.
    \item \textbf{Inference complexity}: $\mathcal{O}(n)$ per token, $\mathcal{O}(r_{kv} \cdot n)$ latent state; path chosen by hardware.
    \item \textbf{Trainable}: trained once; the two decoding paths reuse the same weights.
    \item \textbf{Strengths}: Hardware-adaptive decode with no quality loss; algebraic equivalence guarantees consistency.
    \item \textbf{Limitations}: Same latent-rank tradeoffs as MLA; path selection is a deployment-time heuristic.
\end{itemize}

\subsection{Multi-Head Low-Rank Attention (MLRA)}
Multi-Head Low-Rank Attention (MLRA) \citep{liu_mlra_2026} refactors the single MLA latent into $L$ disjoint sub-heads each of rank $r/L$, so that the KV cache can be partitioned across $L$ devices for 4-way tensor-parallel decoding (each device loads only $1/L$ of the cache). Per-sub-head up-projections reconstruct the corresponding slice of the full heads and are concatenated, yielding a 2.8$\times$ decode speedup over MLA while preserving the global latent budget.

\subsubsection{Mathematical Formulation}
\[
c_{KV} = [c_1, c_2, \ldots, c_L], \quad c_i \in \mathbb{R}^{r/L}, \qquad c_{KV} = W_{DKV}\, x.
\]
\[
k_i = W_{UK}^{(i)}\, c_i, \quad v_i = W_{UV}^{(i)}\, c_i, \qquad
k = [k_1; \ldots; k_L], \quad v = [v_1; \ldots; v_L].
\]
\[
\text{attn} = \texttt{softmax}\!\left(\frac{q\, k^\top}{\sqrt{d_h}}\right) v.
\]

\subsubsection{Algorithmic Pseudocode}
\begin{algorithmic}[1]
\State \textbf{Input:} $x$, sub-head count $L$, per-sub-head projections $W_{UK}^{(i)}, W_{UV}^{(i)}$
\State $c_{KV} \gets W_{DKV}\, x = [c_1, \ldots, c_L]$ \Comment{partition latent into $L$ blocks}
\For{$i = 1, \ldots, L$ (parallel across $L$ devices)}
    \State $k_i \gets W_{UK}^{(i)}\, c_i$, \quad $v_i \gets W_{UV}^{(i)}\, c_i$
\EndFor
\State $k \gets \texttt{concat}(k_1, \ldots, k_L)$, \quad $v \gets \texttt{concat}(v_1, \ldots, v_L)$
\State $\text{weights} \gets \texttt{softmax}(q\, k^\top / \sqrt{d_h})$
\State \textbf{Output:} $\mathbf{Y} = \text{weights} \cdot v$
\end{algorithmic}

\subsubsection{Key Characteristics}
\begin{itemize}[leftmargin=*]
    \item \textbf{Training complexity}: $\mathcal{O}(n^2 \cdot d)$ (full attention on concatenated heads).
    \item \textbf{Inference complexity}: $\mathcal{O}(n)$ per token, $\mathcal{O}(r_{kv} \cdot n)$ state partitioned across $L$ devices.
    \item \textbf{Trainable}: fully trainable; $L$ controls partition granularity.
    \item \textbf{Strengths}: 2.8$\times$ decode speedup; clean 4-way tensor-parallel cache sharding; preserves total latent budget.
    \item \textbf{Limitations}: Requires $L$ to divide $r_{kv}$; more projection matrices than MLA; inter-device all-reduce on output.
\end{itemize}

\subsection{Tucker Attention}
Tucker Attention \citep{klein_tucker_2026} unifies MHA, GQA and MLA under a single Tucker-style factorisation of the query/key/value weight tensors with independent ranks $q_{\text{rank}}, k_{\text{rank}}, v_{\text{rank}}$. MHA is recovered when all ranks equal the hidden size, GQA when $k_{\text{rank}} = v_{\text{rank}} < d$, and MLA when $k_{\text{rank}} = v_{\text{rank}}$ equals the latent rank. Because the factorisation is a pure linear reparameterisation, the method is fully compatible with FlashAttention and RoPE, and yields an order of magnitude fewer parameters for comparable metrics.

\subsubsection{Mathematical Formulation}
\[
\mathbf{Q} = W_{Q,\text{core}}\, W_{Q,\text{factor}}\, x, \qquad
\mathbf{K} = W_{K,\text{core}}\, W_{K,\text{factor}}\, x, \qquad
\mathbf{V} = W_{V,\text{core}}\, W_{V,\text{factor}}\, x.
\]
\[
\text{attn} = \texttt{softmax}\!\left(\frac{\mathbf{Q}\, \mathbf{K}^\top}{\sqrt{d_h}}\right) \mathbf{V}.
\]
Special cases: MHA $\Leftrightarrow q_{\text{rank}} = k_{\text{rank}} = v_{\text{rank}} = d$; GQA $\Leftrightarrow k_{\text{rank}} = v_{\text{rank}} < d$; MLA $\Leftrightarrow k_{\text{rank}} = v_{\text{rank}} = r_{kv}$.

\subsubsection{Algorithmic Pseudocode}
\begin{algorithmic}[1]
\State \textbf{Input:} $x$, ranks $q_{\text{rank}}, k_{\text{rank}}, v_{\text{rank}}$, core/factor matrices
\State $\mathbf{Q} \gets W_{Q,\text{core}}\, W_{Q,\text{factor}}\, x$ \Comment{rank-$q_{\text{rank}}$ factorised query}
\State $\mathbf{K} \gets W_{K,\text{core}}\, W_{K,\text{factor}}\, x$ \Comment{rank-$k_{\text{rank}}$ factorised key}
\State $\mathbf{V} \gets W_{V,\text{core}}\, W_{V,\text{factor}}\, x$ \Comment{rank-$v_{\text{rank}}$ factorised value}
\State \textbf{Apply RoPE to $\mathbf{Q}, \mathbf{K}$ if positional encoding is used}
\State $\text{weights} \gets \texttt{softmax}(\mathbf{Q}\, \mathbf{K}^\top / \sqrt{d_h})$ \Comment{FlashAttention-compatible}
\State \textbf{Output:} $\mathbf{Y} = \text{weights} \cdot \mathbf{V}$
\end{algorithmic}

\subsubsection{Key Characteristics}
\begin{itemize}[leftmargin=*]
    \item \textbf{Training complexity}: $\mathcal{O}(n^2 \cdot d)$ on the reconstructed tensors.
    \item \textbf{Inference complexity}: $\mathcal{O}(n)$ per token; cache size governed by $k_{\text{rank}}$.
    \item \textbf{Trainable}: ranks are hyperparameters; the factorisation is a drop-in weight reparameterisation.
    \item \textbf{Strengths}: Unifies MHA/GQA/MLA; order-of-magnitude parameter reduction at matched metrics; FlashAttention/RoPE compatible.
    \item \textbf{Limitations}: Two matmuls per projection add latency; rank selection is a per-task tuning problem.
\end{itemize}

\subsection{Interleaved Head Attention (IHA)}
Interleaved Head Attention (IHA) \citep{duvvuri_iha_2026} constructs $P$ pseudo-heads per original head (typically $P = H$) where each pseudo query/key/value is a learned linear combination of all $H$ original heads. This induces up to $P^2$ distinct attention patterns per head at an $\mathcal{O}(H^2 P)$ overhead. Theoretically, IHA needs only $\Theta(\sqrt{k}\, n^2)$ parameters versus $\Theta(k\, n^2)$ for MHA on the Polynomial task; empirically it delivers $+10$--$20\%$ on RULER multi-key retrieval, $+5.8\%$ on GSM8K and $+2.8\%$ on MATH-500.

\subsubsection{Mathematical Formulation}
\[
q_{\text{pseudo}}[p] = \sum_{h=1}^{H} \text{mix}_q[p, h]\, q[h], \quad
k_{\text{pseudo}}[p] = \sum_{h} \text{mix}_k[p, h]\, k[h], \quad
v_{\text{pseudo}}[p] = \sum_{h} \text{mix}_v[p, h]\, v[h].
\]
\[
\text{attn}_p = \texttt{softmax}\!\left(\frac{q_{\text{pseudo}}[p]\, k_{\text{pseudo}}[p]^\top}{\sqrt{d_h}}\right) v_{\text{pseudo}}[p],
\qquad
\mathbf{Y}[h] = \frac{1}{P} \sum_{p} \text{attn}_p[h].
\]

\subsubsection{Algorithmic Pseudocode}
\begin{algorithmic}[1]
\State \textbf{Input:} per-head $q[h], k[h], v[h]$, mixing matrices $\text{mix}_q, \text{mix}_k, \text{mix}_v \in \mathbb{R}^{P \times H}$
\For{$p = 1, \ldots, P$}
    \State $q_p \gets \sum_h \text{mix}_q[p,h]\, q[h]$ \Comment{learned mix across heads}
    \State $k_p \gets \sum_h \text{mix}_k[p,h]\, k[h]$, \quad $v_p \gets \sum_h \text{mix}_v[p,h]\, v[h]$
    \State $\text{attn}_p \gets \texttt{softmax}(q_p\, k_p^\top / \sqrt{d_h})\, v_p$
\EndFor
\State \textbf{Output:} $\mathbf{Y}[h] \gets \frac{1}{P} \sum_p \text{attn}_p[h]$ \Comment{average pseudo-head outputs per head}
\end{algorithmic}

\subsubsection{Key Characteristics}
\begin{itemize}[leftmargin=*]
    \item \textbf{Training complexity}: $\mathcal{O}(n^2 \cdot H \cdot P)$ with $P \approx H$, i.e.\ $\mathcal{O}(n^2 \cdot H^2)$.
    \item \textbf{Inference complexity}: $\mathcal{O}(n)$ per token with KV cache; extra $\mathcal{O}(H^2 P)$ mixing overhead.
    \item \textbf{Trainable}: mixing matrices are learned; $P$ is a hyperparameter.
    \item \textbf{Strengths}: $\Theta(\sqrt{k})$ parameter efficiency on Polynomial; large gains on retrieval and reasoning benchmarks.
    \item \textbf{Limitations}: Quadratic-in-heads mixing cost; more attention patterns to compute; added latency at large $H$.
\end{itemize}

\subsection{Grouped-head laTenT Attention (GTA)}
Grouped-head laTenT Attention (GTA) \citep{sun_gta_2025} combines two compressions: (1) a \emph{shared attention map} --- one attention-score tensor per group of heads, reused across all heads in the group, shrinking the key cache; and (2) a \emph{nonlinear value decoder} --- the value cache is compressed to a latent by a down-projection and reconstructed by a SiLU-gated up-projection before the output projection. GTA cuts attention FLOPs by up to 62.5\% versus GQA, shrinks the KV cache by up to 70\%, and yields a 2$\times$ end-to-end inference speedup.

\subsubsection{Mathematical Formulation}
\[
q_g = \texttt{mean}(q[\text{group } g]), \quad k_g = \texttt{mean}(k[\text{group } g]), \quad
\text{attn}_g = \texttt{softmax}\!\left(\frac{q_g\, k_g^\top}{\sqrt{d_h}}\right) \text{ (reused across the group)}.
\]
\[
v_{\text{latent}} = W_{DV}\, x, \qquad v = \text{SiLU}(W_{UV}\, v_{\text{latent}}), \qquad
\mathbf{Y} = \text{attn}_g \cdot v.
\]

\subsubsection{Algorithmic Pseudocode}
\begin{algorithmic}[1]
\State \textbf{Input:} $q, k$ per head, hidden $x$, group partition $\{g\}$
\State $v_{\text{latent}} \gets W_{DV}\, x$ \Comment{compress values to latent (cached)}
\State $v \gets \text{SiLU}(W_{UV}\, v_{\text{latent}})$ \Comment{nonlinear value decode}
\For{each group $g$}
    \State $q_g \gets \texttt{mean}(q[g])$, \quad $k_g \gets \texttt{mean}(k[g])$ \Comment{shared per-group queries/keys}
    \State $\text{attn}_g \gets \texttt{softmax}(q_g\, k_g^\top / \sqrt{d_h})$
    \State $\mathbf{Y}[g] \gets \text{attn}_g \cdot v[g]$ \Comment{one map reused for all heads in $g$}
\EndFor
\State \textbf{Output:} $\mathbf{Y}$
\end{algorithmic}

\subsubsection{Key Characteristics}
\begin{itemize}[leftmargin=*]
    \item \textbf{Training complexity}: $\mathcal{O}(n^2 \cdot G \cdot d_h)$ where $G$ is the number of groups (vs $H$ heads for GQA).
    \item \textbf{Inference complexity}: $\mathcal{O}(n)$ per token; value cache $\mathcal{O}(r_v \cdot n)$, key cache reduced by group sharing.
    \item \textbf{Trainable}: group assignment and latent rank are trainable/hyperparameters.
    \item \textbf{Strengths}: up to 62.5\% attention-FLOP reduction, 70\% KV-cache shrink, 2$\times$ end-to-end speedup.
    \item \textbf{Limitations}: Shared map reduces per-head diversity; nonlinear decoder adds latency; grouping must be chosen carefully.
\end{itemize}

\subsection{Multi-head Temporal Latent Attention (MTLA)}
Multi-head Temporal Latent Attention (MTLA) \citep{deng_mtla_2025} extends MLA along the temporal axis: a hyper-network dynamically merges $m$ temporally adjacent KV-cache vectors into a single slot, shrinking the temporal length by a factor of $m$. A stride-aware causal mask keeps parallel training consistent with autoregressive inference, yielding a 5.3$\times$ decode speedup and 8.3$\times$ GPU-memory reduction versus MHA on En--De speech translation.

\subsubsection{Mathematical Formulation}
\[
c_{KV}^{(t)} = W_{DKV}\, x^{(t)} \in \mathbb{R}^{r_{kv}} \text{ (per-token latent)}.
\]
\[
\bar c_{KV}^{(j)} = \sum_{i=0}^{m-1} \alpha_i\, c_{KV}^{(s j + i)}, \qquad \alpha_i = \texttt{sigmoid}(\text{HyperNet}(c_{KV}^{(s j + i)})),
\]
where $s$ is the stride and $m$ the merge window. Queries attend over merged slots $\bar c_{KV}$ with a stride-aware causal mask $M_{\text{stride}}$:
\[
\text{attn} = \texttt{softmax}\!\left(\frac{q\, \bar k^\top}{\sqrt{d_h}} + M_{\text{stride}}\right) \bar v, \quad \bar k = W_{UK}\, \bar c_{KV},\ \bar v = W_{UV}\, \bar c_{KV}.
\]

\subsubsection{Algorithmic Pseudocode}
\begin{algorithmic}[1]
\State \textbf{Input:} per-token latents $c_{KV}^{(t)}$, merge window $m$, stride $s$
\State \textbf{Merge phase (hyper-network):}
\For{each merged slot $j = 0, 1, \ldots$}
    \State $\alpha_i \gets \texttt{sigmoid}(\text{HyperNet}(c_{KV}^{(s j + i)}))$ for $i = 0, \ldots, m-1$
    \State $\bar c_{KV}^{(j)} \gets \sum_{i=0}^{m-1} \alpha_i\, c_{KV}^{(s j + i)}$ \Comment{gated average over window}
\EndFor
\State $\bar k \gets W_{UK}\, \bar c_{KV}$, \quad $\bar v \gets W_{UV}\, \bar c_{KV}$
\State $\text{weights} \gets \texttt{softmax}(q\, \bar k^\top / \sqrt{d_h} + M_{\text{stride}})$ \Comment{stride-aware causal mask}
\State \textbf{Output:} $\mathbf{Y} = \text{weights} \cdot \bar v$
\end{algorithmic}

\subsubsection{Key Characteristics}
\begin{itemize}[leftmargin=*]
    \item \textbf{Training complexity}: $\mathcal{O}(n^2 \cdot d / m)$ on the merged sequence (length $\lceil n/m \rceil$) plus hyper-network cost.
    \item \textbf{Inference complexity}: $\mathcal{O}(n/m)$ per token over merged slots; $\mathcal{O}(r_{kv} \cdot n / m)$ state.
    \item \textbf{Trainable}: hyper-network and merge window $m$ are learned/chosen; stride-aware mask enforces consistency.
    \item \textbf{Strengths}: 5.3$\times$ decode speedup, 8.3$\times$ GPU-memory reduction vs MHA; temporal compression orthogonal to rank compression.
    \item \textbf{Limitations}: Merge window can blur fine temporal detail; hyper-network adds compute; mask design is non-trivial.
\end{itemize}

\subsection{Architectural Comparison and Synthesis}

\small
\begin{longtable}{p{2.0cm}p{1.8cm}p{1.6cm}p{1.6cm}p{4.5cm}}
\toprule
\textbf{Architecture} & \textbf{Train} & \textbf{Infer} & \textbf{State} & \textbf{Key Characteristic} \\
\midrule
\endhead
Standard Attention & $O(n^2 d)$ & $O(n)$ cache & $O(nd)$ & Perfect expressiveness, full token routing, KV bottleneck. \\
GQA & $O(n^2 d)$ & $O(n)$ cache & $O(h_k \cdot n \cdot d_h)$ & Tunable KV compression, proven at scale, $G \times$ cache reduction. \\
Sigmoid Attention & $O(n^2 d)$ & $O(n)$ cache & $O(nd)$ & Element-wise gating, hardware-efficient, stable at scale with hybrid-norm. \\
RetNet & $O(n^2 d)$ & $O(1)$ & $O(d^2)$ & Multi-scale decay, triple computation mode, fixed forgetting pattern. \\
Mamba & $O(nd)$ & $O(1)$ & $O(d)$ & Selective input-dependent dynamics, linear training and inference, hardware scan. \\
ODE Transformer & $O(kn^2 d)$ & $O(kn)$ & $O(nd)$ & Numerical integration refinement, weight sharing via stages, higher accuracy. \\
Titans & $O(c^2 + nf)$ & $O(c^2 + 1)$ & $O(1)$ & Test-time memory adaptation, extreme context, on-inference parameter updates. \\
MLA & $O(n^2 d)$ & $O(n)$ cache & $O(r_{kv} n)$ & Latent KV of rank $r_{kv}\!\approx\!d/2$, halves cache, matches MHA quality. \\
GQLA & $O(n^2 d)$ & $O(n)$ cache & $O(r_{kv} n)$ & Same MLA weights, two algebraically equivalent hardware-adaptive decode paths. \\
MLRA & $O(n^2 d)$ & $O(n)$ cache & $O(r_{kv} n)$ partitioned & Latent split into $L$ sub-heads for 4-way tensor-parallel decode, 2.8$\times$ speedup. \\
Tucker & $O(n^2 d)$ & $O(n)$ cache & $O(k_{\text{rank}} n)$ & Tucker factorisation unifies MHA/GQA/MLA, order-of-magnitude fewer parameters. \\
IHA & $O(n^2 H^2)$ & $O(n)$ cache & $O(H n d_h)$ & Pseudo-heads as learned head mixes, $\Theta(\sqrt{k})$ params, +10--20\% RULER retrieval. \\
GTA & $O(n^2 G d_h)$ & $O(n)$ cache & $O(r_v n)$ & Shared per-group map + SiLU value decoder, 62.5\% FLOP cut, 70\% cache shrink. \\
MTLA & $O(n^2 d / m)$ & $O(n/m)$ & $O(r_{kv} n/m)$ & Temporal merge of $m$ KV slots, 5.3$\times$ decode, 8.3$\times$ memory reduction. \\
\bottomrule
\end{longtable}
\normalsize

The field exhibits a clear progression along two axes: (i) \textbf{computational efficiency}, moving from $O(n^2)$ to $O(n)$ or $O(1)$, and (ii) \textbf{memory adaptivity}, shifting from static weights to dynamic, test-time updated models. Standard attention remains the expressiveness baseline; Mamba and RetNet represent the practical efficiency frontier; Titans introduces an orthogonal innovation axis (test-time learning). The choice of architecture reflects the fundamental engineering constraint: expressiveness versus deployment cost.

